\documentclass{article}
\usepackage{amsmath,amssymb,amsthm}
\newtheorem{proposition}{Proposition}
\title{Semantic instances of the retained fibre-effect construction}
\date{15 November 2026}\begin{document}\maketitle
\paragraph{Separation.} These are interpretations of the general categorical construction, not its definitions or required application domains. No control protocol, program language or new historical novelty is claimed. The following elementary checks establish the effect-algebra hypotheses; they do not automatically supply a particular application's fibre presentations, body squares or descent certificates.

\section{Deterministic selection}
For $T=\mathrm{Id}$, an effect on a raw fibre is a selected witness. Bind is function application and Fubini is ordinary pairing. The monad is affine, commutative and sends the empty set to the empty set. A native conditional map must send the selected witness to the selected target witness. Raw carriers may retain additional witnesses, but the identity monad supplies no separate probability or canonical possibility interpretation for them. All conditional effects must still be supplied at their declared indices.

\section{Total nondeterministic choice}
Let $T(X)$ be the nonempty subsets of $X$ (or the finite nonempty subsets for the finitary version). The unit is singleton, bind is union over the chosen subset, and Fubini is Cartesian product. Union associativity and distributivity of Cartesian product over unions give the monad and Fubini/bind interchange laws. There is exactly one nonempty subset of a singleton, while the empty set has no nonempty subset. Hence both versions satisfy the affine, commutative and nonempty-fibre hypotheses.

Native conditional transport says that the image of the supplied chosen subset is exactly the target chosen subset. It does not imply surjectivity on the entire raw carrier if that carrier contains elements outside the supplied subset. Ordinary powerset with the empty choice is not affine: it has two effects on a singleton. It is therefore not an unqualified instance of the normalized construction.

\section{A shared external environment}
Fix a nonempty set $R$ and put $T(X)=X^R$. The unit is constant function, and
\[
 (u\mathbin{\gg=}k)(r)=k(u(r))(r),\qquad
 (u\otimes_Tv)(r)=(u(r),v(r)).
\]
Substitution verifies the monad laws and both orders of Fubini agree pointwise. There is one map $R\to1$ and no map $R\to\varnothing$. Thus this reader effect is another nonprobabilistic instance.

The two factors use the \emph{same} environment $r$. This is not independent sampling of two private environments. A supplied base witness and all local effects are evaluated in that common environment; conditional equations hold pointwise in $r$. The general Fubini map therefore does not mean statistical independence across every semantic interpretation. In this instance it means coherent use of an explicit shared parameter.

\section{Finite normalized distributions}
Let $T(X)$ be finitely supported probability distributions. Unit is a point mass, bind is finite weighted summation, and Fubini is the product distribution. The monad laws and interchange follow by reindexing finite sums and multiplying coefficients. Total mass one gives affinity; there is no probability distribution on the empty set.

Here the specified Fubini lift is a chosen product coupling \emph{conditional on the shared base witness}. It is not a uniqueness theorem for all couplings with those marginals. A correlated coupling that does not factor this way is not silently included: its requisite correlation must be placed in the shared witness or represented by a different supplied operation. Zero-weight raw witnesses are retained by the presentation data independently of the distribution support. No conditioning is synthesized from a marginal and no infinite-path measure is supplied by the finitary monad.

\section{Consequences and nonconsequences}
Each instance inherits the finite serial/synchronized formulae only after it has the required supplied fibre effects and body typing. The effective-quotient criterion is tested on those conditional effects, not just on endpoint behavior. In every instance intermediate witnesses must remain retained unless a separate descent condition justifies forgetting them. None of these checks settles arbitrary feedback, partiality, unbounded stochastic limit laws or the whole earlier polynomial independent-interaction criterion.
\end{document}
